Finally, observe that @@M0@@, so the morphisms @@M1@@ and @@M2@@ are obtained by passage to the quotient from the morphisms @@M3@@ and @@M4@@ of 3.7. Thus we have @@M5@@.

By VI 5.8, @@M6@@ for every @@M7@@, since @@M8@@, with @@M9@@ of rank @@M10@@. We therefore obtain a homomorphism:
@@M11@@
and an analogous homomorphism after tensoring with @@M12@@. The homomorphism @@M13@@ is therefore obtained from @@M14@@ (3.7) by passage to the associated graded. Since @@M15@@ is a left inverse of @@M16@@, @@M17@@ is a left inverse of @@M18@@, proving at the same time the injectivity of @@M19@@.

To show that @@M20@@, we must show that for @@M21@@, and @@M22@@, satisfying:
\sgaStableTarget{sga6.e07.eq.4-6-1}@@M23@@
there exists @@M24@@ such that:
@@M25@@

Put:
@@M26@@
Then by VI 5.8, @@M27@@. Next observe that @@M28@@ preserves the filtration (using results of Exposé VIII, which causes no difficulty, since 4.8 is not used to prove this fact). Indeed, it follows from 1.8 ii) that @@M29@@ is a complete-intersection morphism of virtual relative dimension zero (VIII 4.2). Since @@M30@@ is projective, the result follows from VIII 3.2. Applying @@M31@@ to the relation
